\documentclass[11pt]{article}
\usepackage[margin=1in]{geometry}
\usepackage{booktabs}
\usepackage{amsmath}
\usepackage{siunitx}
\usepackage{caption}
\usepackage{hyperref}
\hypersetup{colorlinks=true, linkcolor=blue, urlcolor=blue}

\title{Weighting Scheme: Inverse Volatility (Naive Risk Parity)}
\author{MetaFVG Mean-Reversion Portfolio -- 14 Instruments}
\date{\today}

\begin{document}
\maketitle

\section{Scheme Description}

The portfolio blends 14 fundamentally diversified instruments (7 FX majors/crosses,
4 equity indices, and 3 additional FX pairs spanning different central banks) trading
a Spearman-gated MetaFVG mean-reversion strategy. The current production blend uses
equal notional weights (\(1/14\) each). This report tests whether an
\textbf{inverse-volatility} (naive risk-parity) weighting scheme improves realized
diversification relative to equal weighting.

Each instrument's own standalone daily-return volatility \(\sigma_i\) is computed over
its own non-NaN trading days (days with no closed trade for that instrument are
excluded from that instrument's own vol estimate). Weights are set inversely
proportional to \(\sigma_i\), so instruments with lower standalone volatility receive
higher weight -- equalizing each instrument's naive volatility contribution while
completely ignoring the correlation structure between instruments:

\begin{equation}
w_i = \frac{1/\sigma_i}{\sum_{j=1}^{14} 1/\sigma_j}, \qquad
\sigma_i = \operatorname{std}\big(\{r_{i,t} : r_{i,t} \neq \text{NaN}\}\big)
\end{equation}

On any given day, only instruments with a non-NaN return that day contribute; the
scheme's weights are renormalized to sum to 1 over the subset of instruments trading
that day before the weighted sum is taken. Days on which \emph{all} 14 instruments are
NaN (no trade closed anywhere) are dropped from the blended series.

\section{Weights}

\begin{table}[h]
\centering
\caption{Inverse-volatility weights, sorted descending. Ratio column is
\(w_i / (1/14)\), i.e.\ the multiple of the equal-weight allocation.}
\begin{tabular}{lS[table-format=1.4]S[table-format=1.4]S[table-format=1.6]}
\toprule
Symbol & {Weight} & {Ratio to $1/14$} & {$\sigma_i$ (daily)} \\
\midrule
USDSGD & 0.1351 & 1.8917 & 0.000020 \\
USDCAD & 0.1120 & 1.5679 & 0.000024 \\
EURUSD & 0.1050 & 1.4705 & 0.000026 \\
GBPUSD & 0.0890 & 1.2467 & 0.000030 \\
EURNOK & 0.0839 & 1.1750 & 0.000032 \\
USDCHF & 0.0834 & 1.1670 & 0.000032 \\
USDJPY & 0.0709 & 0.9930 & 0.000038 \\
AUDUSD & 0.0654 & 0.9157 & 0.000041 \\
NZDUSD & 0.0571 & 0.7996 & 0.000047 \\
USDZAR & 0.0495 & 0.6937 & 0.000054 \\
GER40  & 0.0450 & 0.6304 & 0.000060 \\
US500  & 0.0431 & 0.6037 & 0.000063 \\
HK50   & 0.0311 & 0.4357 & 0.000087 \\
JP225  & 0.0293 & 0.4095 & 0.000092 \\
\bottomrule
\end{tabular}
\end{table}

Low-volatility FX pairs (USDSGD, USDCAD, EURUSD, GBPUSD) receive up to 1.9x the
equal-weight allocation, while the higher-volatility equity indices (JP225, HK50)
and equity/EM instruments are cut to roughly 0.4-0.6x equal weight.

\section{Correlation Impact}

\begin{table}[h]
\centering
\caption{Weighted average pairwise correlation and diversification ratio,
equal-weight baseline vs.\ inverse-volatility scheme.}
\begin{tabular}{lS[table-format=1.6]S[table-format=1.6]}
\toprule
Metric & {Equal Weight} & {Inverse Vol} \\
\midrule
Weighted avg.\ pairwise correlation & 0.015615 & 0.015947 \\
Diversification ratio (DR)          & 2.080922 & 1.808789 \\
\bottomrule
\end{tabular}
\end{table}

The diversification ratio is defined as \(\mathrm{DR} = \dfrac{\sum_i w_i \sigma_i}{\sigma_{\text{portfolio}}}\),
where \(\sigma_{\text{portfolio}}\) is the standard deviation of the final blended
daily-return series. Higher DR indicates more realized diversification benefit from
combining the instruments.

\section{Results}

\begin{table}[h]
\centering
\caption{Quantstats risk/return metrics, equal-weight baseline vs.\ inverse-volatility scheme.}
\begin{tabular}{lS[table-format=-1.4]S[table-format=-1.4]}
\toprule
Metric & {Equal Weight} & {Inverse Vol} \\
\midrule
Sharpe                  & 0.5179  & 0.3214 \\
Sortino                 & 0.8170  & 0.5037 \\
Calmar                  & 0.3430  & 0.2001 \\
CAGR                    & 0.0002  & 0.0001 \\
Max drawdown (\%)       & -5.3252 & -5.3252 \\
Volatility (\%, ann.)   & 3.5276  & 3.3172 \\
Tail ratio              & 1.2387  & 1.1027 \\
Ulcer index             & 0.0002  & 0.0002 \\
Kelly criterion         & 0.0668  & 0.0471 \\
VaR (\%)                & -0.0036 & -0.0034 \\
CVaR (\%)                & -0.0072 & -0.0076 \\
Skew                    & 1.5448  & 1.9403 \\
Recovery factor         & 2.2601  & 1.3186 \\
\midrule
Win rate                & 0.3089  & 0.3016 \\
Best day                & 0.000291 & 0.000291 \\
Worst day               & -0.000239 & -0.000239 \\
Total return            & 0.001204 & 0.000702 \\
\bottomrule
\end{tabular}
\end{table}

\section{Takeaway}

The inverse-volatility scheme did not improve realized diversification: the weighted
average pairwise correlation was essentially unchanged (0.01595 vs.\ 0.01561, marginally
\emph{worse}), and the diversification ratio actually fell from 2.08 to 1.81. This makes
sense in hindsight -- inverse-vol weighting ignores the correlation matrix entirely and
simply concentrates capital in the lowest-volatility instruments (USDSGD, USDCAD, EURUSD),
which does nothing to reduce co-movement and instead reduces the effective breadth of the
book. Sharpe dropped substantially, from 0.518 to 0.321 (-38\%), along with Sortino,
Calmar, and recovery factor all deteriorating in tandem. Equal weighting remains the
better scheme of the two tested here; naive risk parity is not a good fit for a portfolio
that was already assembled for fundamental (not volatility-based) diversification.

\end{document}
